\documentclass[11pt]{article}

\usepackage[margin=1in]{geometry}
\usepackage{amsmath,amssymb,mathtools,booktabs,microtype}
\usepackage{enumitem}
\usepackage{xcolor}
\usepackage[numbers,sort&compress]{natbib}
\usepackage[colorlinks=true,linkcolor=blue!50!black]{hyperref}

\title{Direction A: Richer Process Models for Kalman Associative Memory}
\author{}
\date{July 30, 2026}

\newcommand{\R}{\mathbb{R}}
\newcommand{\C}{\mathbb{C}}
\newcommand{\ve}[1]{\boldsymbol{#1}}
\newcommand{\tr}{\intercal}
\newcommand{\I}{\ve{I}}
\newcommand{\N}{\mathcal{N}}
\DeclareMathOperator{\diag}{diag}
\DeclareMathOperator{\blockdiag}{blockdiag}
\DeclareMathOperator{\RePart}{Re}

\begin{document}
\maketitle

\section{Question and scope}

Kalman associative memory separates the evolution of memory from the correction made by a
new key--value observation. Its linear--Gaussian model is
\begin{align}
  \ve{S}_t^\star
    &= \ve{D}_t\ve{S}_{t-1}^\star+\ve{W}_t,
  &
  \ve{W}_t&\sim\N(0,\ve{\Omega}_t),
  \label{eq:base-process}\\
  \ve{v}_t
    &= \ve{S}_t^{\star\tr}\ve{k}_t+\ve{e}_t,
  &
  \ve{e}_t&\sim\N(0,r_t\I).
  \label{eq:base-observation}
\end{align}
The Kalman prediction is
\begin{equation}
  \widehat{\ve{S}}_t=\ve{D}_t\ve{S}_{t-1},
  \qquad
  \widehat{\ve{P}}_t
  =\ve{D}_t\ve{P}_{t-1}\ve{D}_t^{\tr}+\ve{\Omega}_t.
  \label{eq:base-prediction}
\end{equation}
The usual discussion interprets $\ve{D}_t$ as a forgetting operator. More precisely,
vanilla DeltaNet has $\ve{D}_t=\I$ and does not impose global geometric forgetting;
later delta writes can still damage an old binding when their keys overlap it. Explicit
geometric decay enters Gated DeltaNet through $\ve{D}_t=\alpha_t\I$ and KDA through
$\ve{D}_t=\diag(\ve{\alpha}_t)$.

Diagnostics from the current GDN-2 model make the question concrete. Its median learned
amplitude half-life is seven tokens, while only $16.9\%$ of channels exceed 512 tokens and
$7.8\%$ exceed 4096 tokens. At length 8192, the mean effective rank of the memory is
$18.1$ out of 128. These measurements motivate a persistent subspace, but do not by
themselves prove that decay is the only cause: low-dimensional key and write geometry can
also produce low state rank. The hypothesis here is therefore not merely ``add slower
channels.'' It is to give surviving information a controllable latent geometry and route
selected content through it.

The Kalman view permits a broader question. Instead of only asking how quickly a stored
association should shrink, can the process model move it into a latent subspace that has
little effect on current predictions, preserve it there, and make it visible to a later
query? A useful richer process must control three properties separately:
\begin{center}
\begin{tabular}{lll}
\toprule
Property & Controlled by & Meaning \\
\midrule
survival & transition-product norm & how long state amplitude remains \\
phase & rotation angle & how state moves among latent coordinates \\
visibility & write/read routing & whether a query currently observes it \\
\bottomrule
\end{tabular}
\end{center}
For the normal rotation blocks developed below, the transition-product norm is exactly the
product of their radii. Spectral radius alone would not control finite-time survival for a
general non-normal or time-varying transition. Positive diagonal decay already permits
visibility changes through a changing query; what it lacks is a process mechanism that moves
surviving information between latent directions. This note develops a structured linear
$\ve{D}_t$ with that mechanism. A nonlinear
state-dependent process $f_t(\ve{S}_{t-1})$ is possible in an extended-filter formulation,
but it gives up the exact Kalman recursion and most of the current scan structure. The first
step should therefore remain linear in the recurrent state, while allowing its coefficients
to be token dependent and known before the memory scan.


\section{Geometric decay cannot hide by state motion}

Consider the contribution of one innovation written at time $\tau$,
\begin{equation}
  \ve{U}_\tau=\ve{\kappa}_\tau\ve{\delta}_\tau^{\tr},
  \qquad
  \ve{\delta}_\tau
  =\ve{v}_\tau-\widehat{\ve{S}}_\tau^{\tr}\ve{k}_\tau.
\end{equation}
To isolate the process model, first suppress all intervening measurement corrections. The
hypothetical process-only contribution at time $t>\tau$ is
\begin{equation}
  \widetilde{\ve{S}}^{\mathrm{proc}}_{t\leftarrow\tau}
  =\ve{\Phi}_{t:\tau}\ve{\kappa}_\tau\ve{\delta}_\tau^{\tr},
  \qquad
  \ve{\Phi}_{t:\tau}
  =\ve{D}_t\ve{D}_{t-1}\cdots\ve{D}_{\tau+1}.
  \label{eq:process-contribution}
\end{equation}
Applying a query to this process-only state gives
\begin{equation}
  \widetilde{\ve{o}}^{\mathrm{proc}}_{t\leftarrow\tau}
  =\ve{\delta}_\tau
   \underbrace{\ve{\kappa}_\tau^{\tr}
   \ve{\Phi}_{t:\tau}^{\tr}\ve{q}_t}_{\text{visibility of write $\tau$ at time $t$}}.
  \label{eq:visibility}
\end{equation}
Equation~\eqref{eq:visibility} distinguishes process survival from process-only visibility. The norm
$\|\ve{\Phi}_{t:\tau}\ve{\kappa}_\tau\|$ measures how much of the write survives, whereas
its alignment with $\ve{q}_t$ measures how much is currently observable.

For KDA,
\begin{equation}
  \ve{\Phi}_{t:\tau}
  =\diag\!\left(
    \prod_{s=\tau+1}^{t}\alpha_{s,1},\ldots,
    \prod_{s=\tau+1}^{t}\alpha_{s,d_k}
  \right).
  \label{eq:diagonal-product}
\end{equation}
When $0\leq\alpha_{s,i}\leq1$, the magnitude of every scalar mode is nonincreasing.
A fast mode becomes exponentially small and cannot be restored by the process alone. A slow
mode survives, but it is not hidden by the process: it remains in the same coordinate and is
visible whenever the query overlaps that coordinate. Different $\alpha_i$ therefore form a
bank of decay timescales, not a bank of Fourier frequencies. A fixed basis rotation
$\ve{U}\diag(\ve{\alpha})\ve{U}^{\tr}$ still has only positive real eigenvalues and the same
zero-frequency exponential modes in a different coordinate system.

The exact posterior-mean recurrence includes every intervening residual correction. Expanding
the Kalman update gives
\begin{equation}
  \ve{S}_t
  =\underbrace{(\I-\ve{\kappa}_t\ve{k}_t^{\tr})\ve{D}_t}_{\ve{A}_t}
   \ve{S}_{t-1}
   +\ve{\kappa}_t\ve{v}_t^{\tr}.
  \label{eq:closed-loop-carry}
\end{equation}
Define the closed-loop propagation operator
\begin{equation}
  \ve{\Psi}_{t:\tau}=\ve{A}_t\ve{A}_{t-1}\cdots\ve{A}_{\tau+1}.
  \label{eq:closed-loop-product}
\end{equation}
The actual contribution of write $\tau$ to the posterior output at time $t$ is therefore
\begin{equation}
  \ve{o}_{t\leftarrow\tau}
  =\ve{\delta}_\tau
   \ve{\kappa}_\tau^{\tr}\ve{\Psi}_{t:\tau}^{\tr}\ve{q}_t.
  \label{eq:closed-loop-contribution}
\end{equation}
The process-only response in \eqref{eq:process-contribution} isolates what the choice of
$\ve{D}_t$ contributes; \eqref{eq:closed-loop-contribution} is the retrieval quantity the
full layer must preserve. For hidden modes to survive, ordinary measurement keys must weakly
observe those modes, or later residual corrections will erase and mix them.


\section{A damped rotational process model}

General tunable phase requires two-dimensional real planes; a scalar negative pole is the
special Nyquist case $\theta=\pi$. Define the real rotation
\begin{equation}
  \ve{R}(\theta)
  =
  \begin{bmatrix}
    \cos\theta&-\sin\theta\\
    \sin\theta& \cos\theta
  \end{bmatrix}.
\end{equation}
Partition the memory rows into $M$ fixed mode groups. Mode $m$ contains two quadratures of an
$r_m$-dimensional content address and uses
\begin{equation}
  \ve{D}_t^{(m)}
  =\rho_{t,m}\ve{R}(\theta_{t,m})\otimes\I_{r_m},
  \qquad
  0<\rho_{t,m}\leq1,
  \label{eq:oscillator-block}
\end{equation}
with
\begin{equation}
  \ve{D}_t
  =\blockdiag\!\left(
    \ve{D}_t^{(1)},\ldots,\ve{D}_t^{(M)}
  \right).
  \label{eq:spectral-transition}
\end{equation}
An optional fixed orthogonal basis $\ve{U}$ gives
$\ve{U}\ve{D}_t\ve{U}^{\tr}$; in a neural layer this basis can usually be absorbed into the
key and query projections rather than applied to the recurrent state at every token.

Because all rotations within one fixed plane commute, the transition product remains simple:
\begin{equation}
  \ve{\Phi}_{t:\tau}^{(m)}
  =
  \left(\prod_{s=\tau+1}^{t}\rho_{s,m}\right)
  \ve{R}\!\left(\sum_{s=\tau+1}^{t}\theta_{s,m}\right)
  \otimes\I_{r_m}.
  \label{eq:oscillator-product}
\end{equation}
For constant $\rho_m$ and $\theta_m$, the eigenvalues are the complex-conjugate pair
$\rho_m e^{\pm\mathrm{i}\theta_m}$ in the closed unit disk. The homogeneous dynamics are
asymptotically stable for $\rho_m<1$ and only marginally stable for $\rho_m=1$. The radius
$\rho_m$ sets lifetime. For $\rho_m<1$,
the amplitude half-life is
\begin{equation}
  L_{1/2,m}=\frac{\log 2}{-\log\rho_m},
\end{equation}
while $\rho_m=1$ gives an infinite half-life. The angle $\theta_m$ sets frequency and phase
evolution. These are independent parameters.

For a scalar content direction, let $\ve{b}_m\in\R^2$ be its write phase and
$\ve{c}_m\in\R^2$ its read phase. The process impulse response is
\begin{equation}
  h_m(\ell)
  =\ve{c}_m^{\tr}(\rho_m\ve{R}(\theta_m))^\ell\ve{b}_m
  =\rho_m^\ell
   \left(A_m\cos(\ell\theta_m)+B_m\sin(\ell\theta_m)\right)
  \label{eq:damped-mode-response}
\end{equation}
for scalar coefficients $A_m,B_m$ determined by the write and read phases. A bank of such modes
produces
\begin{equation}
  h(\ell)
  =\sum_{m=1}^{M}\rho_m^\ell
   \left(A_m\cos(\ell\theta_m)+B_m\sin(\ell\theta_m)\right),
  \label{eq:filter-bank-response}
\end{equation}
a complex-pole filter bank. It is BIBO stable when every radius is strictly below one;
unit-radius modes are non-expansive without input but accumulate persistent writes. This is
Fourier-like, but it is not generally a Fourier transform of the entire history: damping
makes it an online rational filter, and token-dependent angles make it nonstationary. Complex
phase with decay is already an established sequence-model construction
\citep{sun2023retentive}; negative real eigenvalues are also known to improve linear-RNN state
tracking \citep{grazzi2025statetracking}. The question here is specifically how these modes
interact with Kalman associative writes, uncertainty, and content-addressed retrieval.

The distinction between frequency and timescale is important. A low-frequency aligned mode
has $\cos\theta\approx1$ and can be maximally visible at the next token. A high-frequency mode
can alternate or cancel, but need not be short lived. A desired recent-versus-delayed split
must be designed jointly through radius, frequency, and relative write/read phase.


\section{Hiding in one quadrature and revealing in the other}

Let $\ve{e}_{\mathrm{vis}}=(1,0)^{\tr}$ and
$\ve{e}_{\mathrm{hid}}=(0,1)^{\tr}$. An aligned immediate path has
\begin{equation}
  \ve{b}=\ve{c}=\ve{e}_{\mathrm{vis}},
  \qquad
  h_{\mathrm{direct}}(\ell)
  =\rho^\ell\cos(\ell\theta).
\end{equation}
In particular, $\theta=0$ recovers ordinary geometric decay. A cross-quadrature path has
\begin{equation}
  \ve{b}=\ve{e}_{\mathrm{hid}},
  \qquad
  \ve{c}=\ve{e}_{\mathrm{vis}},
  \qquad
  h_{\mathrm{delay}}(\ell)
  =-\rho^\ell\sin(\ell\theta).
  \label{eq:quadrature-delay}
\end{equation}
It is exactly invisible at insertion and grows as rotation moves it toward the visible axis.
For $0<\rho<1$ and $0<\theta<\pi$, treating $\ell$ as continuous, the first-lobe peak in
magnitude is
\begin{equation}
  \ell^\star
  =\frac{1}{\theta}
   \arctan\!\left(\frac{\theta}{-\log\rho}\right),
  \label{eq:delayed-peak}
\end{equation}
which approaches $\pi/(2\theta)$ as $\rho\to1$. The actual token lag is the maximizing nearby
integer, so the continuous expression can be inaccurate near the Nyquist frequency. A slowly
rotating channel is useful as a delayed channel only if its radius is large enough for the
write to survive until this peak.

This is hidden state relative to the layer readout: it is currently unseen, not erased. For a
sequence of read vectors $\ve{g}_{t+j}$, define the finite-horizon retrieval-observability
Gramian
\begin{equation}
  \ve{W}_{o}(t,L)
  =\sum_{j=0}^{L}
   \ve{\Phi}_{t+j:t}^{\tr}\ve{g}_{t+j}\ve{g}_{t+j}^{\tr}
   \ve{\Phi}_{t+j:t}.
  \label{eq:observability-gramian}
\end{equation}
A direction $\ve{z}$ can be hidden now because $\ve{g}_t^{\tr}\ve{z}=0$, yet remain
recoverable if $\ve{z}^{\tr}\ve{W}_{o}(t,L)\ve{z}>0$. For one oscillator read through
$\ve{e}_{\mathrm{vis}}$, the two-step observability matrix is
\begin{equation}
  \begin{bmatrix}
    \ve{e}_{\mathrm{vis}}^{\tr}\\
    \ve{e}_{\mathrm{vis}}^{\tr}\rho\ve{R}(\theta)
  \end{bmatrix}
  =
  \begin{bmatrix}
    1&0\\
    \rho\cos\theta&-\rho\sin\theta
  \end{bmatrix}.
\end{equation}
It has full rank whenever $\rho\sin\theta\neq0$. Thus the latent quadrature is recoverable,
although very small $\theta$ makes recovery poorly conditioned. The same rotation also causes
periodic echoes; damping or a mixture of frequencies is needed when one delayed peak, rather
than repeated resurfacing, is desired. This Gramian concerns the output read $\ve{g}_t$;
detectability of the Kalman covariance is instead governed by the measurement feature used in
the update.


\section{A Fourier bank gives modulo-age selection}

The frequency analogy can be made exact for age classes by spending state dimension on a
temporal code. This section is an illustrative augmented-state, broadcast-gain mean
recurrence, not yet the exact Kalman update derived in the next section: it changes the state
dimension, write map, and lag-conditioned read in addition to changing $\ve{D}_t$.
Let $\ve{Z}_{m,t}\in\C^{d_a\times d_v}$ be mode $m$, choose
$\omega_m=2\pi m/M$, and broadcast the same associative innovation
$\ve{U}_t=\bar{\ve{\kappa}}_t\ve{\delta}_t^{\tr}$ to every mode:
\begin{equation}
  \ve{Z}_{m,t}
  =\rho e^{\mathrm{i}\omega_m}\ve{Z}_{m,t-1}+\ve{U}_t,
  \qquad m=0,\ldots,M-1.
  \label{eq:complex-delay-write}
\end{equation}
Given a content query $\bar{\ve{q}}_t$ and a requested lag $\ell$, read with the conjugate
phase
\begin{equation}
  \ve{o}_t(\ell)
  =\RePart\!\left[
    \frac{1}{M}\sum_{m=0}^{M-1}
    e^{-\mathrm{i}\omega_m\ell}\ve{Z}_{m,t}^{\tr}\bar{\ve{q}}_t
  \right].
  \label{eq:complex-delay-read}
\end{equation}
For integer age $n=t-\tau$ and requested integer lag $\ell$, one write contributes the phase factor
\begin{equation}
  \frac{1}{M}\sum_{m=0}^{M-1}
  e^{\mathrm{i}\omega_m(n-\ell)}
  =
  \begin{cases}
    1,&n\equiv\ell\pmod M,\\
    0,&\text{otherwise}.
  \end{cases}
  \label{eq:fourier-delay-identity}
\end{equation}
Consequently,
\begin{equation}
  \ve{o}_{t\leftarrow\tau}(\ell)
  =\rho^n
   (\bar{\ve{\kappa}}_\tau^{\tr}\bar{\ve{q}}_t)
   \ve{\delta}_\tau
\end{equation}
when $n\equiv\ell\pmod M$, and it cancels at the other represented age classes. Unrolling all
writes makes the aliasing explicit:
\begin{equation}
  \ve{o}_t(\ell)
  =\sum_{\substack{n\geq0\\n\equiv\ell\ (\mathrm{mod}\ M)}}
   \rho^n\ve{U}_{t-n}^{\tr}\bar{\ve{q}}_t.
  \label{eq:modulo-age-accumulator}
\end{equation}
This is a modulo-$M$ age-addressed accumulator, not an overwriting circular buffer. It selects
one unique age only while the represented history is shorter than $M$; afterward it sums all
aliases, geometrically weighted when $\rho<1$. Content similarity chooses \emph{which}
association to retrieve, while phase matching chooses \emph{which age class} to expose.
Complex arithmetic is not required: conjugate mode pairs correspond to the real $2\times2$
rotation blocks in \eqref{eq:oscillator-block}; the DC mode, and the Nyquist mode for even
$M$, are one-dimensional real exceptions.

A true length-$M$ sliding Fourier delay line would have to remove the expired write,
\begin{equation}
  \ve{Z}_{m,t}
  =\rho e^{\mathrm{i}\omega_m}\ve{Z}_{m,t-1}
   +\ve{U}_t-\rho^M\ve{U}_{t-M}.
  \label{eq:true-fourier-delay-line}
\end{equation}
The expired write can be supplied by an external delay buffer, but it is also recoverable from
the complete modal state. If
$\ve{Z}_{m,t-1}=\sum_{n=0}^{M-1}\rho^n e^{\mathrm{i}\omega_m n}\ve{U}_{t-1-n}$, then
\begin{equation}
  \ve{U}_{t-M}
  =\rho^{-(M-1)}\frac{1}{M}\sum_{j=0}^{M-1}
   e^{\mathrm{i}\omega_j}\ve{Z}_{j,t-1}.
\end{equation}
Substitution into \eqref{eq:true-fourier-delay-line} gives the state-only update
\begin{equation}
  \ve{Z}_{m,t}
  =\rho e^{\mathrm{i}\omega_m}\ve{Z}_{m,t-1}+\ve{U}_t
   -\frac{\rho}{M}\sum_{j=0}^{M-1}
    e^{\mathrm{i}\omega_j}\ve{Z}_{j,t-1}.
  \label{eq:coupled-fourier-overwrite}
\end{equation}
Thus exact overwrite remains fixed dimensional, but it is no longer an independent
block-diagonal oscillator recurrence: the removal is a rank-one coupling across all modes.
That coupling changes both the scan and the covariance structure developed for
\eqref{eq:spectral-transition}.

This construction answers the existence question, but it exposes the cost. It replicates a
content address across $M$ temporal modes, aliases ages separated by $M$, and requires the
read controller to supply a lag or compatible phase. Unequal radii prevent the radial factor
from leaving the Fourier sum and turn the exact age-class response into a soft kernel that
depends on both age and requested lag. Learned complex weights can instead
produce a general query-dependent response
\begin{equation}
  h_t(n)
  =\RePart\!\left[
    \sum_m w_{t,m}\rho_m^n e^{\mathrm{i}n\omega_m}
  \right],
  \label{eq:learned-lag-kernel}
\end{equation}
trading exact age-class selection for a lower-dimensional, task-adaptive temporal filter.


\section{The corresponding Kalman model}

A spectral state may use more rows than the original content key. It is therefore useful to
separate the measurement address from the external key and the read address from the external
query. Let
\begin{equation}
  \ve{S}_t\in\R^{d_s\times d_v},
  \qquad
  \ve{P}_t\in\R^{d_s\times d_s},
  \qquad
  \ve{E}_t,\ve{F}_t\in\R^{d_s\times d_k},
\end{equation}
and define
\begin{equation}
  \ve{h}_t=\ve{E}_t\ve{k}_t\in\R^{d_s},
  \qquad
  \ve{g}_t=\ve{F}_t\ve{q}_t\in\R^{d_s},
\end{equation}
where $\ve{E}_t$ chooses the measured modes and phases and $\ve{F}_t$ chooses the modes and
phases visible to the query. Under a block-isotropic covariance the measurement lift also
selects the corresponding write modes; under a dense covariance, cross-correlations can rotate
the Kalman gain elsewhere. The generalized associative-memory filter is
\begin{align}
  \widehat{\ve{S}}_t
    &=\ve{D}_t\ve{S}_{t-1},
  &
  \widehat{\ve{P}}_t
    &=\ve{D}_t\ve{P}_{t-1}\ve{D}_t^{\tr}+\ve{\Omega}_t,
  \label{eq:spectral-kalman-predict}\\
  \ve{\kappa}_t
    &=\frac{\widehat{\ve{P}}_t\ve{h}_t}
      {r_t+\ve{h}_t^{\tr}\widehat{\ve{P}}_t\ve{h}_t},
  &
  \ve{S}_t
    &=\widehat{\ve{S}}_t+\ve{\kappa}_t
      (\ve{v}_t-\widehat{\ve{S}}_t^{\tr}\ve{h}_t)^{\tr},
  \label{eq:spectral-kalman-update}\\
  \ve{P}_t
    &=(\I-\ve{\kappa}_t\ve{h}_t^{\tr})\widehat{\ve{P}}_t,
  &
  \ve{o}_t&=\ve{S}_t^{\tr}\ve{g}_t.
  \label{eq:spectral-kalman-read}
\end{align}
This is not a new filtering rule. It is the same Kalman filter with a richer latent geometry.
The original model is recovered by $d_s=d_k$, $\ve{E}_t=\ve{F}_t=\I$, and a diagonal
$\ve{D}_t$. The explicit $\ve{E}_t,\ve{F}_t$ make one limitation precise: $\ve{D}_t$ can
preserve and move information, but it cannot by itself define what ``hidden from this query''
means.

The broadcast write in the Fourier example is not automatically the gain produced here.
It is Kalman-consistent only under a symmetric measurement lift and covariance that yield
identical gain blocks, or it can be used as a fixed-gain approximation. Its role above is to
prove an age-selection property of an augmented spectral state, not to claim that $\ve{D}_t$
alone supplies a delay memory.

Rotational dynamics have a particularly convenient covariance approximation. For mode $m$,
let its dimension be $d_m=2r_m$. Assume the full covariance and process noise are block
diagonal and isotropic inside every block,
\begin{equation}
  \ve{P}_{t-1}
  =\blockdiag_m(p_{t-1,m}\I_{d_m}),
  \qquad
  \ve{\Omega}_t
  =\blockdiag_m(\nu_{t,m}\I_{d_m}).
\end{equation}
Then the prediction is scalar per mode:
\begin{equation}
  \widehat{\ve{P}}_t^{(m)}
  =\left(\rho_{t,m}^2p_{t-1,m}+\nu_{t,m}\right)\I_{d_m},
  \qquad
  \widehat p_{t,m}=\rho_{t,m}^2p_{t-1,m}+\nu_{t,m}.
  \label{eq:block-isotropic-predict}
\end{equation}
The rotation cancels from covariance prediction while still rotating the memory mean. The
resulting gain is
\begin{equation}
  \ve{\kappa}_t^{(m)}
  =\frac{\widehat p_{t,m}\ve{h}_t^{(m)}}
  {r_t+\sum_j\widehat p_{t,j}\|\ve{h}_t^{(j)}\|_2^2}.
  \label{eq:block-isotropic-gain}
\end{equation}
An exact measurement update leaves the block-isotropic family: it creates rank-one anisotropy
inside every observed block as well as correlations across blocks. This is the same closure
problem that motivates projecting the current diagonal-information KDN rather than calling a
diagonal posterior exact. A trace-matching projection back to the block-isotropic family would
give
\begin{equation}
  p_{t,m}
  =\widehat p_{t,m}
   -\frac{\widehat p_{t,m}^2\|\ve{h}_t^{(m)}\|_2^2}
    {d_m\left(r_t+\sum_j\widehat p_{t,j}
    \|\ve{h}_t^{(j)}\|_2^2\right)}.
  \label{eq:block-isotropic-project}
\end{equation}
This projection is an approximation, not an exact block Kalman posterior. It is also a
sequential coupled recurrence because its denominator depends on every block, so it does not
give the independent M\"obius scan used by the current diagonal KDN.

For a scan-compatible first model, use a calibrated scalar information state
$c_{t,m}=1/p_{t,m}$ per mode. Prediction is exact within the block-isotropic family,
\begin{equation}
  \widehat c_{t,m}
  =\frac{c_{t-1,m}}
   {\rho_{t,m}^2+\nu_{t,m}c_{t-1,m}},
  \qquad
  \widehat p_{t,m}=\frac{1}{\widehat c_{t,m}}.
  \label{eq:block-info-predict}
\end{equation}
The measurement update is projected and standardized as
\begin{equation}
  u_{t,m}
  =\frac{d_s}{d_m r_t}\|\ve{h}_t^{(m)}\|_2^2,
  \qquad
  c_{t,m}=\widehat c_{t,m}+u_{t,m}.
  \label{eq:block-info-update}
\end{equation}
The literal trace projection of the information increment would omit the factor $d_s$.
Equation~\eqref{eq:block-info-update} instead applies the same coordinate standardization as
the current diagonal KDN: for an $\ell_2$-normalized feature spread across $d_s$ coordinates,
$\mathbb{E}[\|\ve{h}_t^{(m)}\|_2^2]\approx d_m/d_s$, so each block receives an expected $O(1/r_t)$
increment. This is a calibrated approximation, not an exact posterior projection.

Combining \eqref{eq:block-info-predict} and \eqref{eq:block-info-update} gives the per-mode
M\"obius recurrence
\begin{equation}
  c_{t,m}
  =\frac{(1+u_{t,m}\nu_{t,m})c_{t-1,m}
          +u_{t,m}\rho_{t,m}^2}
         {\nu_{t,m}c_{t-1,m}+\rho_{t,m}^2},
  \label{eq:block-info-mobius}
\end{equation}
represented by the associative matrix
\begin{equation}
  \ve{M}_{t,m}
  =
  \begin{bmatrix}
    1+u_{t,m}\nu_{t,m} & u_{t,m}\rho_{t,m}^2\\
    \nu_{t,m} & \rho_{t,m}^2
  \end{bmatrix}.
  \label{eq:block-info-matrix}
\end{equation}
The gains in \eqref{eq:block-isotropic-gain} are then available from an independent prefix
scan for each mode before the mean-memory scan. This is the uncertainty approximation assumed
by the recommended spectral KDN below.


\section{Scheduled resurfacing versus recall on cue}

A fixed oscillator makes a memory visible according to elapsed time. That is useful for
delay-sensitive prediction, but it does not by itself retrieve an arbitrary fact when a
semantic cue arrives at an arbitrary future time. There are three increasingly expressive
ways to add cue control.

\paragraph{Query-controlled observability.}
Let the query projection emit both quadratures and a mode gate. At ordinary tokens it can
suppress persistent modes; at a recall cue it can read the matching hidden phase. In
\eqref{eq:spectral-kalman-read}, this is a token-dependent $\ve{F}_t$ and requires no
state-dependent transition. It is the simplest and most scan-friendly mechanism. The usual
learned query projection can already implement part of this behavior, but an explicit
per-mode gate makes the intended separation testable.

\paragraph{Input-conditioned phase increments.}
Use
\begin{equation}
  \rho_{t,m}=\exp(-\Delta_{t,m}\lambda_m),
  \qquad
  \theta_{t,m}=\Delta_{t,m}\omega_m+\gamma_{t,m},
  \label{eq:dynamic-pole-parameterization}
\end{equation}
where $\Delta_{t,m}>0$ is an input-dependent step size, $\lambda_m\geq0$ is a decay rate,
$\omega_m$ is a base frequency, and $\gamma_{t,m}$ is a bounded input-dependent phase kick.
A retrieval cue can rotate a latent mode toward the visible phase before the current read.
Because the coefficients depend only on input-known quantities, log-radii and angles still
compose by prefix addition. The caveat is that one phase kick rotates every association in
that mode; content selection must still come from the key/query match.

\paragraph{Active/archive routing.}
Allocate paired active and archive subspaces. A write gate routes durable information toward
the archive quadrature, and a read gate chooses when the archive participates in the output.
Continuous $\rho\ve{R}(\theta)\otimes\I$ gives scheduled, periodic exchange between active
and archive coordinates; it does not keep an item archived until a cue. Cue-controlled archive
semantics require either $\theta_t=0$ between explicit transfer events or a query gate that
fully suppresses the archive until retrieval. The lifts $\ve{E}_t$ and $\ve{F}_t$ then provide
content-conditioned storage and recall. This gives the clearest semantics for ``hide and
later query,'' at the cost of either a larger state or fewer content dimensions per subspace
under a fixed state budget.

For arbitrary cue-driven recall, query-controlled observability should be present before
dynamic phase kicks are added. Otherwise the model may learn a clock that solves fixed-delay
training examples but fails when query times and recall order are randomized.


\section{Parallelism and implementation constraints}

The proposed transition remains structured, but it is not a drop-in use of the current
KDA/GDN-2 delta-family kernels.
\begin{enumerate}[leftmargin=1.7em,itemsep=4pt]
  \item Applying block rotations to $\ve{S}_t$ costs $O(d_s d_v)$, the same order as diagonal
  decay, with a larger constant for paired multiply-adds. A dense learned change of basis
  would instead cost $O(d_s^2d_v)$ and should be absorbed into projections or avoided.

  \item Fixed-plane transitions compose associatively. Each mode can represent its prefix
  product by cumulative log-radius and cumulative phase, or by complex multiplication. This
  keeps the process prediction scan-friendly.

  \item Once gains are input known, the full mean recurrence is the affine scan
  \begin{equation}
    \ve{S}_t
    =(\I-\ve{\kappa}_t\ve{h}_t^{\tr})\ve{D}_t\ve{S}_{t-1}
     +\ve{\kappa}_t\ve{v}_t^{\tr}.
  \end{equation}
  It is associative in principle, but efficient DeltaNet chunking exploits more structure
  than abstract associativity. The paired/complex key scores and decay normalization require
  a new derivation and kernel; equal asymptotic work does not imply equal throughput.

  \item A block-isotropic uncertainty prediction stays scalar per mode, as shown in
  \eqref{eq:block-isotropic-predict}, and the calibrated update
  \eqref{eq:block-info-mobius} supplies a per-mode gain scan. A full $2\times2$ covariance is
  a useful next step for a scalar-content plane. When $r_m>1$, an exact update additionally
  creates correlations among content directions; cross-mode correlations still require
  projection or dense filtering.

  \item If $\theta_{t,m}$ depends on the sequential memory $\ve{S}_{t-1}$, it is not known
  before the memory scan. Initial models should use fixed angles, token-local angles, or angles
  derived from separately scan-computable prefix statistics.
\end{enumerate}

A conservative development ladder is
\begin{center}
\begin{tabular}{cll}
\toprule
Stage & Transition & Question isolated \\
\midrule
A0 & pair-tied channels, $\theta=0$ & matched no-rotation control \\
A1 & fixed $\theta_m$, learned $\rho_m$ & do complex poles help? \\
A2 & fixed poles, learned write/read phase & can state be hidden then read? \\
A3 & query-dependent read phase & can recall occur on an arbitrary cue? \\
A4 & token-dependent phase increments & does dynamic reorganization help? \\
\bottomrule
\end{tabular}
\end{center}


\section{Limits and failure modes}

\begin{enumerate}[leftmargin=1.7em,itemsep=4pt]
  \item \textbf{No extra capacity for free.} A fixed-dimensional linear state cannot retain an
  unbounded number of independent values losslessly. Oscillation changes temporal coding and
  interference; it does not eliminate superposition capacity limits.

  \item \textbf{Aliasing.} A finite frequency grid confuses ages modulo its period. Damping,
  multiple incommensurate frequencies, resets, or a sufficiently large temporal code can
  soften this, but each choice has a state or optimization cost.

  \item \textbf{Persistent noise.} With $\rho=1$, additive writes and process noise do not
  decay. In a unitary mode that remains unobserved, covariance can grow without bound. Stable
  operation needs $\rho<1$, reliable overwrite/reset behavior, or sufficiently frequent
  observability.

  \item \textbf{Phase mismatch.} Rotating the memory address without a compatible key/query
  lift can make an association inaccessible rather than usefully latent. The transition,
  observation feature, and read feature must be designed together.

  \item \textbf{Periodic leakage.} One oscillator repeatedly reveals a stored value. A mixture
  can localize one lag through interference, but imperfect learned phases may produce periodic
  false recalls.

  \item \textbf{Optimization collapse.} Learned angles may collapse to zero, reducing the
  model to ordinary decay, or cluster at a few phases. Fixed frequency initialization and
  explicit diagnostics are needed before concluding that the spectral model is unused.
\end{enumerate}


\section{Falsifiable experiments}

All comparisons should match total recurrent state dimension, parameter count, and, where
possible, FLOPs. An augmented Fourier bank is otherwise advantaged simply by storing more
numbers.

\paragraph{Lag-resolved associative recall.}
Train MQAR with a broad distribution of write--query gaps and report accuracy by exact lag,
including lags outside the training range. Compare ordinary untied diagonal KDN/KDA, pair-tied
$\theta=0$ modes with exactly the same lifts and query gates as the spectral model, signed-real
$\theta\in\{0,\pi\}$ modes, fixed oscillators, learned write/read phases, and
query-demodulated modes. Periodic accuracy dips would directly reveal aliasing or phase
mismatch.

\paragraph{Survival versus visibility probe.}
Write one tagged association, feed controlled distractors, and record both the process-only
lineage $\ve{\Phi}_{t:\tau}\ve{\kappa}_\tau$ and the actual closed-loop lineage
$\ve{\Psi}_{t:\tau}\ve{\kappa}_\tau$. For each, measure norm and visibility under a fixed
read phase; an intervention difference between runs with and without the tagged write is an
equivalent empirical measurement. Holding the read fixed prevents changing query routing from
creating an apparent delayed peak even for diagonal decay. The defining spectral hypothesis
is that closed-loop lineage norm can remain high while current visibility is low, followed by
a later rise in visibility.

\paragraph{Cue-controlled recall.}
Randomize query times and the order in which stored items are requested. This prevents a fixed
oscillator clock from solving the task. Ablate the temporal read phase, content router, and
input-conditioned phase kick separately. A useful archive must respond to the cue rather than
only to elapsed time.

\paragraph{Collision and capacity tests.}
At fixed lag, vary the number of simultaneously stored key--value pairs and the number of
distractors. This tests whether improved lag selectivity is purchased by a loss of content
capacity. Measure effective state rank and per-mode write/read energy in addition to task
accuracy.

\paragraph{Core ablations and diagnostics.}
The essential controls are $\theta=0$ damping only, signed-real
$\theta\in\{0,\pi\}$, $\rho=1$ rotation only, damped rotation, aligned versus
cross-quadrature routing, fixed versus learned versus token-dependent angles, and diagonal
versus block-isotropic uncertainty. Log learned poles $(\rho_m,\theta_m)$, half-lives, phase
usage, process-only and closed-loop impulse responses, state norms, covariance scale,
throughput, and gradient stability.


\section{Recommended first model}

The first implementation should be a \emph{spectral KDN} rather than a fully general learned
transition:
\begin{enumerate}[leftmargin=1.7em,itemsep=3pt]
  \item pair fixed state planes and use
  $\ve{D}_t^{(m)}=\rho_{t,m}\ve{R}(\theta_m)\otimes\I_{r_m}$;
  \item retain zero-angle modes spanning both fast and slow radii, and include the pair-tied
  zero-angle model with identical lifts and gates as the main no-rotation control;
  \item initialize the remaining angles from log-spaced target delays
  $L_m$ through $\theta_m\approx\pi/(2L_m)$ and give them radii whose lifetimes exceed
  $L_m$;
  \item let key and query projections learn the two quadratures, with an explicit query gate
  over modes;
  \item use block-isotropic process noise and the calibrated M\"obius information scan in
  \eqref{eq:block-info-mobius}, so the gain is available before the mean scan; and
  \item test variable-lag and randomized-cue recall before introducing token-dependent phase
  kicks.
\end{enumerate}

The central claim is deliberately narrow. Positive geometric decay decides how much of a
memory remains. Damped rotations add a separate mechanism that decides where the surviving
memory sits relative to the current readout. With phase-aware write and read maps, this can
hide an association in a latent quadrature and later expose it. A complete Fourier bank gives
exact selection of an age class, but without explicit overwrite it accumulates aliases rather
than becoming a finite delay buffer. Whether the spectral process helps language modeling
depends on closed-loop survival, content routing, cue-controlled observability, interference,
and efficient kernels, all of which the proposed experiments measure directly.

\bibliographystyle{plainnat}
\bibliography{main}

\end{document}
